% ===================================================================
%  図表第 3 号 — 幼年と青年。
%
%  Ceci n'est PAS une transcription : c'est une traduction, faite en
%  2026, en japonais d'aujourd'hui. Voir texto/ja/10-zuhyo-01.tex
%  pour la consigne, qui vaut pour les seize fichiers.
%
%  « 七月 » n'est pas en gras ici, non plus qu'en ido : le fac-simile
%  l'oublie entre « 六月 » et « 八月 », et c'est gravuri/korekti.json
%  qui le retablit, colonne par colonne, sur la page de lecture.
% ===================================================================

%%K t03-apar-1 apar
\VUsaut{3.0mm}
\VUtitre{15.0pt}{60}{図表第 3 号}
\VUsaut{1.6mm}
\VUtitre{11.4pt}{0}{幼年と青年。}
\VUsaut{3.4mm}

%%K t03-apar-2 apar
（第三と第四の図表はこう組んである。四つの\VUgras{家庭の場面}のうちに、
\VUgras{天気}、\VUgras{年齢}、\VUgras{家族}、\VUgras{衣服}、\VUgras{健康}の
基本の語彙が出るように。）
\VUsaut{4.0mm}

%%K t03-tit-1 sub
\VUcentre{12.0pt}{0}{\textit{第一場。}}
\VUsaut{3.0mm}

%%K t03-tit-2 sub
\VUpk{10.2pt}{40}{村の洗礼。}
\VUsaut{2.4mm}

%%K t03-tit-3 sub
\VUpk{10.2pt}{40}{春。}
\VUsaut{3.0mm}

%%K t03-c1-01-1 p
1. --- いまは\VUgras{春}、\VUgras{三月}、\VUgras{四月}あるいは\VUgras{五月}、
\VUgras{一週間}のうちのある日 --- \VUgras{月曜}、\VUgras{火曜}あるいは
\VUgras{水曜}である。朝の\VUgras{十時}。

%%K t03-c1-02-1 p
2. --- \VUgras{天気}はよく、\VUgras{空気}は澄んでいる。日が照っている。小さな
\VUgras{雲} \textsuperscript{(2)} が幾つか青い\VUgras{空} \textsuperscript{(1)} に
のぼってゆく。

%%K t03-c1-03-1 p
3. --- \VUgras{木}にはまだほとんど\VUgras{葉}がない。細い\VUgras{ポプラ}
\textsuperscript{(3)} と高い\VUgras{プラタナス} \textsuperscript{(17)} には、まだ
芽しかついていない。

%%K t03-c1-04-1 p
4. --- \VUgras{つばめ} \textsuperscript{(4)} が帰ってきて、うれしげに鳴きながら
あの\VUgras{尖塔} \textsuperscript{(7)} のまわりを廻っている。それは\VUgras{鐘楼}
\textsuperscript{(5)} の尖塔で、鐘楼は村の\VUgras{教会} \textsuperscript{(6)} の上に
立っている。

%%K t03-tit-4 sub
\VUsaut{3.4mm}
\VUpk{10.2pt}{40}{教会。}
\VUsaut{3.0mm}

%%K t03-c1-05-1 p
5. --- この教会はたいそう質素で、大きな\VUgras{門} \textsuperscript{(13)} と
大きな\VUgras{時計} \textsuperscript{(8)} がある。\VUgras{短い針}
\textsuperscript{(9)} は X の字を指している。それは\VUgras{時}を示す。
\VUgras{長い方} \textsuperscript{(10)} は XII（正午）を指している。それは
\VUgras{分}を示す。

%%K t03-c1-06-1 p
6. --- 鐘楼では\VUgras{鐘} \textsuperscript{(11)} がうれしげに鳴っている。屋根の
一番高いところに石の小さな\VUgras{十字架} \textsuperscript{(12)} が立っている。

%%K t03-c1-07-1 p
7. --- 教会には大きな\VUgras{門} \textsuperscript{(13)} だけでなく、脇に小さな戸
もある。\VUgras{司祭} \textsuperscript{(15)} は\VUgras{長衣}を着て、この戸から
\VUgras{祭服室} \textsuperscript{(14)} を出て、\VUgras{司祭館} \textsuperscript{(19)}
の方へ行くところである。

%%K t03-c1-08-1 p
8. --- そのそばにいるのは\VUgras{牛乳売りの女} \textsuperscript{(24)} と彼女の
\VUgras{ロバ} \textsuperscript{(23)} である。町から帰るところで、子供たちのために
よい牛乳を届けてきたのである。

%%K t03-tit-5 sub
\VUsaut{4.0mm}
\VUpk{10.2pt}{40}{果樹園。}
\VUsaut{3.4mm}

%%K t03-c1-09-1 p
9. --- アリボロン殿（ロバ）はこの朝の道行きにたいそう満足である。
\VUgras{花} \textsuperscript{(20)} の香に満ちた空気を心地よげに吸っている。
その花は隣の\VUgras{果樹園}の\VUgras{果樹} \textsuperscript{(18)} を覆って
いるのである。あの
\VUgras{桃の木} \textsuperscript{(18)} と\VUgras{杏の木} \textsuperscript{(19)} は
すっかり桃色と白で、よい実りを約束している。

%%K t03-c1-10-1 p
10. --- そのあいだ小さな\VUgras{蜜蜂} \textsuperscript{(22)} が\VUgras{巣箱}
\textsuperscript{(21)} からそこへ飛んでゆき、羽音を立てながら甘い\VUgras{蜜}を
集める。

%%K t03-c1-11-1 p
11. --- ところが何ごとであろう。村の子供たちが駆けてきて、おびえた
\VUgras{鵞鳥} \textsuperscript{(25)} を追い散らしている。公証人の子の
\VUgras{洗礼}のあと、行列が教会から出てくるところなのである。

%%K t03-c1-12-1 p
12. --- 小さなジャックは\VUgras{揺籠} \textsuperscript{(26)} の中で、うれしげな
鐘の音に目を覚ます。姉のマドレーヌも行列を見たくて、母に戸口まで来て髪を結い、
着物を着せてくれと頼む。

%%K t03-c1-13-1 p
13. --- 母は承知する。自分の\VUgras{頭巾} \textsuperscript{(28)}、自分の
\VUgras{胴着} \textsuperscript{(29)}、自分の\VUgras{前掛} \textsuperscript{(30)} を
つけて、戸の前の小さな椅子に坐りに来る。マドレーヌは\VUgras{下着}
\textsuperscript{(31)} と\VUgras{胴衣} \textsuperscript{(32)} しかつけていない。

%%K t03-c1-14-1 p
14. --- なんとうれしそうなことか。彼女は大好きな花を忘れている。自分の
\VUgras{なでしこ} \textsuperscript{(34)}、その上には\VUgras{蝶}
\textsuperscript{(35)} が止まっている。自分の\VUgras{チューリップ}
\textsuperscript{(36)}、自分の\VUgras{すずらん} \textsuperscript{(37)}、それは
\VUgras{植木鉢} \textsuperscript{(38)} に入って\VUgras{塀} \textsuperscript{(33)} の
上に並んでいる。それに美しい生垣をなす自分の\VUgras{薔薇} \textsuperscript{(39)}。
彼女は行列の婦人たちの立派な衣裳ばかり眺めている。

%%K t03-c1-15-1 p
15. --- 村の男の子は衣裳をあまり気にかけない。\VUgras{砂糖菓子}
\textsuperscript{(55)} や\VUgras{小銭} \textsuperscript{(52)} をもらおうと、前へ
押し出し、押し合っている。その一人が泣いている \textsuperscript{(40)}。

%%K t03-c1-16-1 p
16. --- もう一人は\VUgras{犬} \textsuperscript{(42)} の足を踏んでしまい、犬は
きゃんきゃん鳴いて逃げ、\VUgras{雌鶏} \textsuperscript{(43)} とその\VUgras{雛}
\textsuperscript{(44)} を追い散らす。

%%K t03-tit-6 sub
\VUsaut{5.0mm}
\VUpk{10.2pt}{40}{洗礼の行列。}
\VUsaut{4.0mm}

%%K t03-c1-17-1 p
17. --- 行列の先頭を歩くのは\VUgras{乳母} \textsuperscript{(45)} である。長い
リボンのついた美しい\VUgras{頭巾} \textsuperscript{(46)} をかぶっている。
\VUgras{生まれたばかりの子} \textsuperscript{(47)} を抱いていて、その子は全身
\VUgras{レース} \textsuperscript{(48)} にくるまれている。

%%K t03-c1-18-1 p
18. --- 乳母の後ろには\VUgras{名づけ親の父} \textsuperscript{(49)} と
\VUgras{名づけ親の母} \textsuperscript{(53)} が来る。二人は砂糖菓子と小銭を撒いて
いる。名づけ親の父は\VUgras{手袋} \textsuperscript{(51)} をはめ、\VUgras{シルク
ハット} \textsuperscript{(50)} をかぶっている。名づけ親の母は手に美しい
\VUgras{花束} \textsuperscript{(54)} を持っている。

%%K t03-c1-19-1 p
19. --- その後ろに子供の\VUgras{伯父} \textsuperscript{(56)} と\VUgras{伯母}
\textsuperscript{(58)} が見える。伯母は洒落た\VUgras{帽子} \textsuperscript{(59)} を
かぶり、伯父は黒い\VUgras{フロックコート} \textsuperscript{(57)} と白い胴衣を着て
いる。次に\VUgras{従兄} \textsuperscript{(60)} と\VUgras{従姉} \textsuperscript{(61)}
が来る。従姉は生まれたばかりの子の\VUgras{姉} \textsuperscript{(62)}、小さな
ベルトの手を引いている。

%%K t03-c1-20-1 p
20. --- ベルトのもう一人の\VUgras{兄} \textsuperscript{(63)} は後ろを歩いている。
一人の\VUgras{親戚} \textsuperscript{(64)} に手を貸している。家の\VUgras{友}
\textsuperscript{(65)} たちがそのあとに続く。

%%K t03-tit-7 sub
\VUsaut{4.6mm}
\VUpk{10.2pt}{40}{見物人。}
\VUsaut{3.4mm}

%%K t03-c1-21-1 p
21. --- 行列のそばに\VUgras{乞食} \textsuperscript{(66)} が一人いて、自分の
\VUgras{松葉杖} \textsuperscript{(67)} にもたれている。施しを乞うている。

%%K t03-c1-22-1 p
22. --- 反対側にはたいそう\VUgras{行儀のよい} \textsuperscript{(68)} 男の子が
いる。帽子を取って、顔見知りに\VUgras{「おはよう」}と言う。

%%K t03-c1-23-1 p
23. --- 村人たちは洗礼の行列が通るのを参照うと家から出てきている。
\VUgras{木綿の帽子}をかぶった\VUgras{村の男} \textsuperscript{(69)} が見え、その
そばに\VUgras{村の女} \textsuperscript{(70)} がいる。次に\VUgras{老婆}
\textsuperscript{(71)} がいて、自分の\VUgras{杖} \textsuperscript{(72)} にもたれて
いる。最後に\VUgras{粉屋} \textsuperscript{(73)} が広い\VUgras{フェルト帽}
\textsuperscript{(74)} をかぶっており、\VUgras{木靴} \textsuperscript{(75)} をはいた
若い村の女が自分の子に歩き方を教えている。

%%K t03-c1-24-1 p
24. --- たいそうにぎやかな眺めである。
\VUsaut{4.0mm}

%%K t03-tit-8 sub
\VUcentre{12.0pt}{0}{\textit{第二場。}}
\VUsaut{3.0mm}

%%K t03-tit-9 sub
\VUpk{10.2pt}{40}{公の祭。}
\VUsaut{2.4mm}

%%K t03-tit-10 sub
\VUpk{10.2pt}{40}{水上の競技と舟の競漕。}
\VUsaut{3.0mm}

%%K t03-c2-01-1 p
1. --- \VUgras{夏}が来た。\VUgras{六月}（七月または\VUgras{八月}）の焼けつく
ような日が、若者たちを水浴びと舟遊びに誘う。

%%K t03-c2-02-1 p
2. --- 川の右岸では漕ぎ手たちが水上の競技と競漕に出ようと着物を脱いでいる。

%%K t03-c2-03-1 p
3. --- その一人は\VUgras{靴} \textsuperscript{(1)} を脱いでいる。紐を解いて
（ボタンを外して）いる。仲間の二人はもう支度ができている。いまは\VUgras{運動着}
\textsuperscript{(2)} と\VUgras{水着} \textsuperscript{(3)} しか身につけていない。
仲間を待ちながら話している。向う岸の市と祭のことを話している。

%%K t03-c2-04-1 p
4. --- そのそばの地面には仲間の\VUgras{着物}が置いてある。一\VUgras{足}の
\VUgras{半長靴} \textsuperscript{(4)}、\VUgras{靴} \textsuperscript{(5)}、
\VUgras{靴下} \textsuperscript{(6)}、\VUgras{フェルト}\textsuperscript{(7)} と
\VUgras{麦藁} \textsuperscript{(8)} の帽子、それに\VUgras{ネクタイ}
\textsuperscript{(9)} が一本。

%%K t03-c2-05-1 p
5. --- 一人の若者が自分の\VUgras{上着} \textsuperscript{(10)} を念入りにたたんだ。
それを手拭の上に置く。隣の者がやがてその手拭で二人の\VUgras{服}を包むのである。

%%K t03-c2-06-1 p
6. --- 六人目の若者は遅れている。頭にはまだ\VUgras{鳥打帽} \textsuperscript{(11)}
をかぶっている。\VUgras{シャツ} \textsuperscript{(12)} も、\VUgras{襟}
\textsuperscript{(13)} も、\VUgras{カフス} \textsuperscript{(14)} も脱いでいない。
手には\VUgras{胴衣} \textsuperscript{(17)} を持ち、まだ\VUgras{ズボン}
\textsuperscript{(16)} をはき、\VUgras{ズボン吊} \textsuperscript{(15)} をかけている。
次の競漕には間に合うまい。

%%K t03-c2-07-1 p
7. --- 若者たちのそばに\VUgras{薊} \textsuperscript{(18)} の生えた小道が水際へ
下りていて、そこでは老ジャックが\VUgras{葦} \textsuperscript{(19)} のあいだで
夕方に放つ\VUgras{花火} \textsuperscript{(20)} を支度している。

%%K t03-tit-11 sub
\VUsaut{4.4mm}
\VUpk{10.2pt}{40}{競漕。}
\VUsaut{3.4mm}

%%K t03-c2-08-1 p
8. --- 川では何人かの婦人が日傘の下に\VUgras{小舟} \textsuperscript{(21)} で
出ている。競漕を見ているのだが、それはたいそう心の躍るものである。どの
\VUgras{競漕艇} \textsuperscript{(22)} にも四人の\VUgras{漕ぎ手} \textsuperscript{(24)}
が力いっぱい漕ぎ、\VUgras{舵手} \textsuperscript{(26)} が\VUgras{舵}
\textsuperscript{(25)} を握って薄い舟を導いている。\VUgras{櫂} \textsuperscript{(23)}
が拍子を合わせて水を打ち、軽い舟が目もくらむ速さで走る。

%%K t03-c2-09-1 p
9. --- 橋脚のそばでは何人かの若者が\VUgras{横滑り竿} \textsuperscript{(28)} の賞を
争っている。その一人がしなってすべる竿の上を用心して進み、先に結んだ小さな
\VUgras{旗} \textsuperscript{(29)} まで届こうとしている。運のわるい\VUgras{相手}
\textsuperscript{(30)} は水に落ちてしまった。泳いで岸へ戻っている。

%%K t03-c2-10-1 p
10. --- 年上の若者たちは\VUgras{水上の試合} \textsuperscript{(32)} をしている。
\VUgras{河岸} \textsuperscript{(33)} のすぐそばである。大勢の\VUgras{人}
\textsuperscript{(31)} が\VUgras{橋} \textsuperscript{(27)} の\VUgras{欄干}にもたれ、
\VUgras{岸}に詰めかけてこれを見ている。しかし何人かの見物人は振り返って
\VUgras{登り竿} \textsuperscript{(45)} の遊びを見、あるいは\VUgras{賞の授与}へ向う
\VUgras{町役所の行列}を眺めている。

%%K t03-c2-11-1 p
11. --- まず\VUgras{わんぱく小僧} \textsuperscript{(35)} が幾人か\VUgras{楽隊}
\textsuperscript{(36)} の前を走っている。次に\VUgras{町長} \textsuperscript{(37)}、
その助役たち、小さな町の\VUgras{町議会}が来る。しんがりを行くのは
\VUgras{消防夫} \textsuperscript{(38)} である。

%%K t03-tit-12 sub
\VUsaut{5.0mm}
\VUpk{10.2pt}{40}{市。}
\VUsaut{3.6mm}

%%K t03-c2-12-1 p
12. --- \VUgras{市の広場} \textsuperscript{(39)} は人であふれている。人々は
いろいろな\VUgras{小屋}を見に来る。\VUgras{木馬} \textsuperscript{(40)}、
\VUgras{曲馬} \textsuperscript{(41)} --- そこではもう演し物が始まっている。それに
\VUgras{公けの舞踏場} \textsuperscript{(42)}、町の若い男女がそこで心ゆくまで
\VUgras{踊り}を楽しむ。

%%K t03-c2-13-1 p
13. --- 奥には\VUgras{闘牛場} \textsuperscript{(44)} が見え、勇ましいスペインの
\VUgras{闘牛士}がそこで怒った\VUgras{牡牛} \textsuperscript{(45)} と戦う。次に
\VUgras{山車滑り} \textsuperscript{(49)} と\VUgras{射的小屋} \textsuperscript{(46)}
があり、人々はそこで\VUgras{銃}を扱い\VUgras{的}を撃つ稽古をする。

%%K t03-c2-14-1 p
14. --- そのそばに\VUgras{市の芝居小屋} \textsuperscript{(47)} があって、
\VUgras{喜劇}と\VUgras{悲劇}を演じる。すぐ隣に\VUgras{巨人} \textsuperscript{(48)}
と\VUgras{小人}の小屋が立っている。前者の\VUgras{背丈}は二メートル十五、後者は
やっと子供ほどの高さしかない。

%%K t03-c2-15-1 p
15. --- 最後に大きな\VUgras{猛獣の小屋} \textsuperscript{(50)} がある。その
\VUgras{野獣}たちのうちには、力強い\VUgras{爪}をもつ\VUgras{虎}
\textsuperscript{(51)}、濃い\VUgras{たてがみ}をもつ\VUgras{獅子}
\textsuperscript{(52)}、二頭の\VUgras{麒麟} \textsuperscript{(53)}、それに
\VUgras{象} \textsuperscript{(54)} がいて、この象は自分の\VUgras{鼻}をじつに巧みに
使う。

%%K t03-c2-16-1 p
16. --- この獣たちを仕込む\VUgras{猛獣使い} \textsuperscript{(55)} が小屋の戸口に
立っている。
\VUsaut{6.0mm}
\VUfilet{22mm}
